% §8.5 What Comes Next
%
% Cross-reference resolution status:
%
% DEFINED here:
%   \label{sec:ch8-what-comes-next}              -- this section
%
% ALREADY DEFINED (chapter 5):
%   \Cref{ch:approximate-sets}                   -- chapter 5
%
% ALREADY DEFINED (chapter 6):
%   \Cref{ch:set-composition}                    -- chapter 6
%
% ALREADY DEFINED (chapter 7):
%   \Cref{ch:classification-measures}            -- chapter 7
%
% ALREADY DEFINED (chapter 8 §§8.1--8.4):
%   \Cref{sec:entropy-floor}                     -- §8.3
%   \Cref{thm:entropy-floor}                     -- §8.3
%   \Cref{sec:bloom-gap}                         -- §8.4
%   \Cref{prop:bloom-overhead}                   -- §8.4
%
% FORWARD REFS (Part III, IV, V stubs from Plan 0):
%   \Cref{ch:approximate-functions}              -- chapter 9
%   \Cref{ch:optimal-bernoulli-sets}             -- chapter 14
%   \Cref{ch:perfect-hash-filters}               -- chapter 15
%   \Cref{ch:not-bernoulli}                      -- chapter 16
%
% §8.5 closes Part II and signposts Parts III, IV, V. Short by design:
% four paragraphs, one per part transition. The chapter has spent its
% mathematical budget on the entropy floor and the Bloom gap; the
% section's job is to name the next destinations and stop.

\section{What Comes Next}
\label{sec:ch8-what-comes-next}

Part~II began with the approximate set as a noisy bit per universe
element, the object \Cref{ch:approximate-sets} built up from the
binary symmetric channel of Part~I. From that atom the part developed
three things: the algebra of approximate sets under union,
intersection, and complement (\Cref{ch:set-composition}); the
estimation discipline for false-positive and false-negative rates and
the derived diagnostic measures that follow from them
(\Cref{ch:classification-measures}); and the information-theoretic
floor that pins the per-element space cost of any FNR-zero
representation (this chapter). The approximate set as an atomic object
is now fully described. Its rates are estimable, its compositions are
predictable, its space cost is bounded below by a number that depends
only on the false-positive rate, and the standard Bloom filter sits a
known $\log_2 e$ above that number. Part~II is done.

Part~III generalizes from sets to maps. \Cref{ch:approximate-functions}
lifts the type $\Bernoulli{\BoolType}$, the noisy bit, to
$\Bernoulli{T}$ for arbitrary algebraic types $T$. The model-order
ceiling of $2$ that the binary channel imposed in Part~I lifts
accordingly: an approximate map into a $k$-element output type is
characterized by a $k \times k$ channel matrix, not a single
$(\fprate, \fnrate)$ pair. Chapters~10 through~13 then develop, in
order, sketches (count-min, hyperloglog, and the family of streaming
estimators), randomized algorithms whose acceptance probability is
itself the object of interest, approximate relations as the lifting
of approximate sets from one column to two, and the formal
algebraic-type framework that ties the type constructors of
$\Bernoulli{T}$ to the error propagation rules Part~II derived for
sets. Maps generalize sets in the same way that channel matrices
generalize the BSC.

Part~IV builds the constructive payoff. \Cref{ch:optimal-bernoulli-sets}
constructs the Bernoulli Hash Function, the approximate set that
achieves the entropy floor of \Cref{thm:entropy-floor} asymptotically.
The construction abandons the flat shared bit array of the standard
Bloom design (the structural choice \Cref{sec:bloom-gap} identified as
responsible for the $\log_2 e$ overhead) in favor of a perfect-hash-based
fingerprint encoding. \Cref{ch:perfect-hash-filters} constructs the
counterpart on the other one-sided axis, the FPR-zero family of
perfect hash filters, which give exact membership at the cost of a
controlled false-negative rate. Together these two chapters close
out the constructive question that Part~II's lower bound leaves open.

Part~V draws the boundary. \Cref{ch:not-bernoulli} catalogues the
systems that are \emph{not} Bernoulli, the cases in which the
independent-error model breaks down: errors that are correlated across
inputs, rates that drift in time, adversarial query distributions that
defeat the uniform-sampling assumption underlying FPR estimation. The
catalogue is itself useful, since it tells the practitioner when the
machinery of Parts~II through~IV applies and when it does not.
Chapters~17 and~18 then exploit the model's edges in a different way,
turning to cipher maps and encrypted search: settings in which the
controlled error of a Bernoulli map serves as a cryptographic resource
rather than a deployment constraint. The thread that runs through all
five parts is the same: an atom of controlled error, its algebra, its
floor, its optimal constructions, and the boundary that fences off
what the atom cannot describe.
